\chapter{Results}\label{chap:results}

The performance and robustness of the proposed Equatorial Plasma Bubble (EPB) detection architecture were evaluated through statistical validation, model interpretability analysis, and continuous temporal application. 

\subsection{Performance Metrics and Cross-Validation}
To make sure the model could generalize and wasn't just memorizing the training data, we ran a 5-fold cross-validation. The entire pipeline, from the 2DPCA feature extraction to the stacking ensemble and SMOTE, performed remarkably consistently across all data partitions. In fact, four out of the five folds hit perfect scores (1.0000) across Accuracy, Sensitivity, Specificity, and F1-Score. The third fold dipped slightly, logging an Accuracy of 0.9902, a Sensitivity of 0.9804, a Specificity of 1.0000, and an F1-Score of 0.9901. Averaging these out, the cross-validation yielded an Accuracy and F1-Score of $99.80\% \pm 0.44\%$. These numbers suggest the system can identify plasma structures with near-perfect reliability on this dataset.

\subsection{Model Interpretability}
We used the SHAP TreeExplainer to investigate the physical relevance of the features driving these decisions. Heatmap analyses showed that when an image actually contained an EPB, the model almost always focused its attention on the upper half of the optical field of view. This makes perfect physical sense, as plasma bubbles typically appear as vertical depletion structures rising into the upper quadrants. The lower half of the images rarely influenced the final decision weight. Interestingly, we also found a very low correlation between the absolute SHAP values and the 2DPCA reconstruction error. This implies that the decision tree ensemble and the mathematical image reconstruction are actually looking at different underlying image properties to confirm a detection.

\subsection{Temporal Stability Analysis}
Beyond static classification, we put the model through a continuous temporal analysis to see how it would hold up during a long operational shift. We fed it a complete, uninterrupted sequence of nocturnal airglow observations from the imager located at Bom Jesus da Lapa - Bahia. The sequence was recorded on July 22, 2025, running from 21:56 to 08:23 UT. Across those 230 consecutive images of a quiet ionosphere (meaning zero EPB events), the model showed strong immunity to false positives. The average EPB probability hovered at just $0.5\%$ for the entire night, and the absolute maximum noise peak only reached $17.7\%$, staying safely below the threshold that would trigger a classification alert. These findings confirm that the model is stable and reliable enough for real-time monitoring.
